% Source: https://artofproblemsolving.com/community/c3h2454554_mathcounts_2019_state_sprint_21
% Stable identity: aops-topic-2454554-q01
\textbf{(b) first.} For a point $X$ chosen uniformly from $[-1,1]$, the probability that it lands in an interval is the interval's length divided by $2$. For every $\varepsilon>0$, the event $X=0$ is contained in the event $-\varepsilon\le X\le\varepsilon$, of probability $\varepsilon$. So $P(X=0)\le\varepsilon$ for every $\varepsilon>0$, and hence $P(X=0)=0$. Events of probability $0$ do not change any probability, so we may ignore the case where a point is exactly at $0$ (whether or not ``between'' would count it).

\textbf{(a)} Let the points be $X$ and $Y$. Ignoring the probability-$0$ event that one of them is $0$, the origin lies between them exactly when one is negative and the other positive. Each point is negative with probability $\dfrac12$ and positive with probability $\dfrac12$, independently, so
\[P=P(X<0,\,Y>0)+P(X>0,\,Y<0)=\dfrac14+\dfrac14=\boxed{\dfrac12}.\]

\begin{asy}
usepackage("amsmath");
size(3.5cm);
defaultpen(fontsize(10pt));
fill((-1,0)--(0,0)--(0,1)--(-1,1)--cycle,rgb(0.85,0.9,1)); fill((0,-1)--(1,-1)--(1,0)--(0,0)--cycle,rgb(0.85,0.9,1));
draw(box((-1,-1),(1,1)),linewidth(0.8));
draw((-1.3,0)--(1.35,0),Arrow(6)); draw((0,-1.3)--(0,1.35),Arrow(6));
label("$X$",(1.35,0),E); label("$Y$",(0,1.35),N);
label("$-1$",(-1,0),SW,fontsize(8pt)); label("$1$",(1,0),SE,fontsize(8pt)); label("$1$",(0,1),NE,fontsize(8pt)); label("$-1$",(0,-1),SE,fontsize(8pt));
\end{asy}

In the square of possible pairs $(X,Y)$, the favourable pairs fill the two shaded quarters.
